% Gerado por tcc/analise/resultados_c3.py -- não edite à mão; rode o script de novo.
\begin{table}[htb]
\caption{Médias anuais por grupo em 31 de dezembro (\%)}\label{tab:anuais}
\centering\footnotesize\setlength{\tabcolsep}{4pt}
\begin{tabular}{lrrrrrr}
\toprule
Ano & \multicolumn{2}{c}{Dívida bruta / Ativo} & \multicolumn{2}{c}{Exigível / Ativo} & \multicolumn{2}{c}{Dívida CP / Dívida total} \\
\midrule
 & F & NF & F & NF & F & NF \\
2018 & 23,9 & 26,4 & 50,2 & 47,8 & 29,2 & 41,4 \\
2019 & 27,0 & 28,3 & 56,1 & 53,3 & 23,0 & 25,4 \\
2020 & 29,6 & 30,5 & 59,0 & 56,0 & 17,2 & 25,8 \\
2021 & 26,0 & 35,4 & 54,5 & 60,1 & 21,9 & 47,0 \\
2022 & 24,9 & 36,1 & 54,9 & 59,0 & 19,6 & 35,9 \\
2023 & 26,2 & 32,8 & 55,5 & 56,8 & 18,9 & 36,8 \\
2024 & 24,7 & 28,2 & 55,6 & 54,6 & 16,1 & 39,1 \\
2025 & 26,4 & 26,9 & 57,2 & 62,1 & 19,1 & 54,3 \\
\bottomrule
\end{tabular}
\fonte{Elaborado pelo autor (2026), com dados da Economática.}
\nota{F: familiares (10 empresas); NF: não familiares (6 empresas; 5 na Dívida CP / Dívida total). Média simples entre as empresas de cada grupo.}
\end{table}
